\newcommand{\ModuleID}{EL-A5}
\newcommand{\ModuleTitle}{Textual Comparison and Advanced Commentary}
\newcommand{\ModuleStage}{Stage IV — Advanced Reading and Composition}
\newcommand{\ModuleUnit}{A5}
\newcommand{\ModuleOutcome}{Preserve a diplomatic base, classify witness differences, build a minimal apparatus, and write concise grammatical commentary}
\newcommand{\ModulePrerequisites}{EL-A4 and access to the exact identified witnesses required by the chosen collation}
\newcommand{\ModuleExtensionPractice}{Worksheets IV.17 and IV.20 remain optional variant-classification and rubrical-register practice}
\newcommand{\ModuleNext}{EL-A6, Controlled Imitation}
\newcommand{\ModuleMastery}{Produce an accurate diplomatic base, a transparent classified apparatus against two named witnesses, and five concise notes whose morphology, syntax, sense, parallel, and alternative are independently checkable.}
\newcommand{\ModuleDelayNote}{\par\medskip\textbf{Scope of the two-session core.} This packet proves collation and commentary. It does not prove independent rubrical-register reading; complete Worksheet IV.20 before claiming that separate skill.}
\newcommand{\ModuleLessonSource}{\CourseRoot/shared/content/volume4/all}
\newcommand{\ModuleExerciseSource}{\CourseRoot/shared/exercises/volume3_4/volume4}
\newcommand{\ModuleWorkedBank}{IV}
\newcommand{\ModuleExerciseBank}{IV}
\newcommand{\ModuleWorksheetVariants}{B,C}
\newcommand{\ModuleScopeNote}{The selected core is collation plus commentary. Variant classification is supplied in transfer practice, while the full rubrical worksheet remains an explicit extension.}
\newcommand{\ModulePractice}{%
  \SubmoduleExtension{A minimal collation}
    {A collation is reproducible only when the base, witnesses, lineation, uncertainties, and normalizations stay visible.}
    {Before Worksheet IV.18, make a five-line mock apparatus with the project-composed base \Latin{Deus pacem donat}; let witness B print \Latin{Deus donat pacem} and witness C print \Latin{Dominus pacem donat}. Classify both without choosing a preferred reading.}
    {State which difference can affect emphasis and which changes the lexeme.}
    {B has transposition, potentially affecting order or emphasis without changing proposition; C has lexical substitution. The exercise is explicitly project-composed and demonstrates notation only, not a received variant.}
  \SubmoduleExtension{The six-line commentary}
    {A concise note becomes checkable when every line answers a different evidentiary question.}
    {Choose one form from the actual base and draft lines for morphology, governor/function, clause relation, contextual sense, checked parallel, and alternative or rejected parse.}
    {Cover the form and reconstruct which evidence belongs on each line.}
    {The model requires one distinct answer per line and an exact witness or course reference for the parallel. A theological reflection cannot replace morphology, and a fabricated alternative is worse than a reason no material alternative survives.}
  \SubmoduleExtension{Rubrical register}
    {Instructional force belongs to discourse context; ordinary present, participial, and gerundive morphology still requires ordinary parsing.}
    {Analyze the lesson's explicitly project-composed rubric-style example, then rewrite its circumstance as a finite clause and state the stylistic change.}
    {Why is present indicative not automatically a special rubrical tense?}
    {The model parses subject, finite present, governed motion, and compressed circumstance normally. The rewrite makes the relation explicit but changes compression; register supplies directive force, not a new tense paradigm.}
  \SubmoduleExtension{Worked collation-to-commentary process}
    {The commentary must be downstream of the established text, not a way to decide the text invisibly.}
    {For one line of Worksheet IV.18, show the complete sequence: diplomatic transcription, differences, classification, normalized aid, selected reading if needed, then six-line note.}
    {Mark the first step at which editorial judgment is permitted.}
    {Transcription and difference recording precede judgment. A substantive decision may be made only after witnesses and classifications are visible; the note must cite that record and preserve residual uncertainty.}
}
